\documentclass[a4paper,12pt]{article}
\usepackage[utf8]{inputenc}
\usepackage[T1]{fontenc}
\usepackage[ngerman]{babel}
\usepackage{amsmath}
\usepackage{amssymb}
\usepackage{enumitem}
\usepackage{geometry}
\geometry{a4paper, margin=2.5cm}
\usepackage{parskip}

\title{Einsendeaufgaben -- Aufgabe 6.5\\[0.5em]
\large Modul 61112: Lineare Algebra}
\author{}
\date{}

\begin{document}
\maketitle

\section*{Aufgabe 6.5}

Wir zeigen zunächst, dass $\chi_f$ in Linearfaktoren zerfällt. Um $\chi_f$ zu
ermitteln, definieren wir eine geeignete Basis. Als ersten Vektor wählen wir ein
beliebiges $v_1 \in \mathrm{Bild}(f) \setminus \{0\}$, dies ist möglich, da
$\mathrm{Bild}(f) \setminus \{0\} \neq \emptyset$ wegen $\mathrm{Rg}(f) = 1$.
$v_1$ lässt sich zu einer Basis $B_{v_1} = (v_1, v_2, \ldots, v_n)$ erweitern.

Da $\mathrm{Rg}(f) = 1$, gilt
\[
  f(v_i) = k_i v_1 \quad \text{mit } k_i \in \mathbb{K} \text{ für alle } v_1, \ldots, v_n.
\]

Demnach ist
\[
  {}_{B_{v_1}}M_{B_{v_1}}(f) = \begin{pmatrix}
    k_1 & k_2 & \cdots & k_n \\
    0 & 0 & \cdots & 0 \\
    \vdots & \vdots & & \vdots \\
    0 & 0 & \cdots & 0
  \end{pmatrix}.
\]

Wir können demnach $\chi_f$ ermitteln
\[
  \chi_f = \det \begin{pmatrix}
    X-k_1 & -k_2 & \cdots & -k_n \\
    0 & X & & \\
    \vdots & & \ddots & \\
    0 & 0 & \cdots & X
  \end{pmatrix} = X^{n-1}(X - k_1).
\]

$\chi_f$ zerfällt also in Linearfaktoren und demnach existiert eine Basis $B$,
sodass ${}_{B}M_{B}(f)$ in jordan'sche Normalform ist.

Nun untersuchen wir noch den möglichen Normalformen für alle
$f \in \mathrm{End}(V)$ mit $\mathrm{Rg}(f) = 1$.

\textbf{\underline{Fall 1: $k_1 = 0$}}

Dann gilt $\chi_f = X^n$ und demnach ist der einzige Eigenwert $0$.
${}_{B}M_{B}(f)$ ist also in nilpotenter Normalform.

\textbf{\underline{Fall 2: $k_1 \neq 0$}}

Da $a(k_1) = 1$ ist $J(k_1, (1)) = (k_1)$. Wenn dann $\dim V > 1$ ist $0$ der
zweite Eigenwert und der Jordanblock von $0$ ist wie in Fall 1 in nilpotenter
Normalform.

\end{document}
